%KerMLformat.tex

\usepackage[parfill]{parskip}    % Activate to begin paragraphs with an empty line rather than an indent
\usepackage{color}
\usepackage{calc}			%need to overcome: ! Missing number, treated as zero. <to be read again>
\usepackage{epstopdf}

%% LaTeX - Article customise

%%% PACKAGES
\usepackage{cite} % will properly break citation lists across lines
\usepackage{booktabs} 	% for much better looking tables
\usepackage{array} 			% for better arrays (eg matrices) in maths
\usepackage{paralist} 		% very flexible & customisable lists (eg. enumerate/itemize, etc.)
\usepackage{verbatim} 		% adds environment for commenting out blocks of text & for better verbatim
%conflicts with tocloft \usepackage{subfigure} 	% make it possible to include more than one captioned figure/table in a single float

%  THEOREMS
\usepackage{amsthm} 		% for theorem styles
\usepackage{thmtools}      % for theorem tools \declaretheorem and \declaretheoremstyle 
\usepackage{thm-restate}  %  restatable theorems
%\usepackage{amsopn}		% for \DeclareMathOperator
\usepackage{soul} 			% for underlining

%  FONTS and SYMBOLS
\usepackage{amssymb} 	% for math symbols
\usepackage{amsmath}		% for \DeclareMathOperator
\usepackage{semantic} 		% for logic symbols and inference rules
\usepackage{txfonts}			% for \llbracket and \rrbracket
\usepackage{MnSymbol}	% for \nsqsubset \ \nmodels and \mathsf
\usepackage{accents}     % for \underaccent
\usepackage{courier}  % make typewriter font Courier so \texttt has bold
%\usepackage{phonetic}      % for \riota
%\usepackage[fleqn]{amsmath} 	% for math words like wp
%\usepackage[mathscr]{eucal} % for Euler math script fonts
%\usepackage{eufrak} 		% for Euler Fraktur font

%  GRAPHICS
\usepackage{graphicx}
\usepackage{tikz}				%graphics package, uses PGF
\usepackage{graphpap}		% for graphpaper in the picture environment
\DeclareGraphicsRule{.tif}{png}{.png}{`convert #1 `dirname #1`/`basename #1 .tif`.png}
%\input xy							% for XY-pic to draw graphs
%\xyoption{all}					% use XY-graph feature "combinatoric drawing language"

\usepackage{ragged2e} 	% for FlushRight and FlushLeft justifications
%\usepackage{graphs}		% for lattice drawing--doesn't work with TeXShop
%\usepackage{calc}				% for expressions in length calculations
\usepackage{ifthen}			% for if-then-else
\usepackage{wrapfig}		% to wrap text around figures with \begin{wrapfigure} environment
\usepackage{makeidx}		%MakeIndex
\usepackage{fancyhdr}     %header and footer style
\usepackage{lastpage}     %determines LastPage number
\usepackage{supertabular}  %breaks-up tables bigger than one page
%hyperref want to be last package
\usepackage[hyperfootnotes=false]{hyperref}		% makes clickable hyperlinks
\usepackage{xr-hyper}    %external references  use \zref{LRM-x} for LRM references

% LISTINGS
\usepackage{listings}      %listings package for BLESS/AADL program text
%\usepackage{minitoc}      %table of contents for each Part
\usepackage{chngcntr}  % prevent reset of paragraph counter
\usepackage{fancyvrb}  % fancy verbatim for BNF
%\usepackage{alltt}

%%%  CHAPTER APPEARANCE
\usepackage[Lenny]{fncychap}
%\usepackage[Conny]{fncychap}
\usepackage{titlesec}
\usepackage[nottoc,notlof,notlot]{tocbibind} % Put the bibliography in the ToC
\usepackage[titles]{tocloft} % Alter the style of the Table of Contents
\usepackage{makeidx}		% for using MakeIndex

% PAGE/MARGIN PARAMETERS
\setlength\textwidth{450 pt}
\setlength\textheight{8.5in}
\setlength\topmargin{0 pt}
\setlength\oddsidemargin{20 pt}
\setlength\evensidemargin{20 pt}
\setlength\headheight{10pt}

%HEADER-FOOTER
\pagestyle{fancy}
\setlength\headheight{23pt}
\lhead{\textsf{\nouppercase{\leftmark}}}
\rhead{\textsf{\nouppercase{\rightmark}}}
%\rhead{\textsf{\thesection}}
%\chead{}
%\rfoot{\color{red}\textsf{\blessversion}} 
%\rfoot{\textsf{\nouppercase{\rightmark}}}
%\cfoot{\multitudelogo{0.04}}
\cfoot{}
\rfoot{-\thepage-}
%\rfoot{\color{red}\textsf{[\blessversion]}}
\renewcommand\headrule{\hrule height 2pt width\headwidth}


%%%%%%%%%% NUMBERING %%%%%%%%%%%%%%%%%%
%\renewcommand{\thechapter}{\arabic{chapter}}
\renewcommand{\thepart}{\Roman{part}}
\renewcommand{\thechapter}{\arabic{chapter}}
\renewcommand{\thesection}{\thechapter.\arabic{section}}
\renewcommand{\thesubsection}{\thesection.\arabic{subsection}}
\renewcommand{\thesubsubsection}{\thesubsection.\roman{subsubsection}}

% COLORS
\definecolor{commentcolor}{gray}{0.6}
\definecolor{codebackgroundcolor}{rgb}{0.9,0.99,0.9}
\definecolor{glossarycolor}{rgb}{1.0,0.0,1.0}
\definecolor{paragraphnumbercolor}{rgb}{0.99,0.60,0.60}
%colors for Metamath proofs
\definecolor{wffcolor}{rgb}{0,0,0.9}
\definecolor{wfflstcolor}{rgb}{0.0,0.5,0.3}
\definecolor{formulacolor}{rgb}{0,0.1,1}
\definecolor{variablecolor}{rgb}{1,0.2,0}
\definecolor{classcolor}{rgb}{0.7,0.2,0.6}
\definecolor{actioncolor}{rgb}{1,0.8,0.2}
\definecolor{assertioncolor}{rgb}{1.0,0.7,0.0}
\definecolor{blcolor}{rgb}{0.1,0,.8}
\definecolor{leancolor}{rgb}{1.0,0,1.0}
\definecolor{leanyellow}{rgb}{0.7,0.7,0}
\definecolor{assertedactioncolor}{rgb}{0,1,0}
\definecolor{setmmcolor}{rgb}{0.2,0.3,0.1}

% MATH COLORING
%well-formed formulas
\newcommand{\mw}[1]{{\color{wffcolor}#1}}
%well-formed formula lists
\newcommand{\mwl}[1]{{\color{wfflstcolor}#1}}
%set variables
\newcommand{\msc}[1]{{\color{variablecolor}#1}}
%classes
\newcommand{\mc}[1]{{\color{classcolor}#1}}
%metamath type-names wff class setvar
\newcommand{\mm}[1]{{\color{gray}#1}}
%assertions
\newcommand{\mass}[1]{{\color{assertioncolor}#1}}
%actions
\newcommand{\mact}[1]{{\color{actioncolor}#1}}
%asserted actions
\newcommand{\maa}[1]{{\color{assertedactioncolor}#1}}


%SYMBOL for empty class
\newcommand{\emptyclass}{\mc{\emptyset}}


%  size of text in listings
\newcommand{\listingsize}{\footnotesize}
%define a language for listings
\lstdefinelanguage{kerml}
	{ 
	keywords={},
	morekeywords=[2]{ 
about, abstract, alias, all, and, as, assoc, behavior, binding, bool, by, chains, class, crosses,
classifier, comment, composite, conjugate, conjugates, conjugation, connector, datatype, 
default, dependency, derived, differences, disjoining, disjoint, doc, else, end, expr, 
false, feature, featured, featuring, filter, first, flow, for, from, function, hastype, if, 
intersects, implies, import, in, inout, interaction, inv, inverse, inverting, istype, 
language, member, metaclass, metadata, multiplicity, namespace, nonunique, not, null, of, 
or, ordered, out, package, portion, predicate, private, protected, public, readonly, 
redefines, redefinition, references, rep, return, specialization, specializes, step, 
struct, subclassifier, subset, subsets, subtype, succession, then, to, true, type, typed, 
typing, unions, var, xor
},
	sensitive=true,
	morecomment=[l]{//},
	columns=[r]fixed,
	basicstyle=\ttfamily\listingsize,
	tabsize=2,
%	otherkeywords= {
%	(, ),  [, ], ;,  ~, @, #, &, ^, |, *, **, +, -, /, ->, ., .., :, :: %, \{, \}, :>, :>>, ::>, <, <=, =, :=, ==, ===, !=, !==, >, >=, ?, ??, .?
%  }
	}
	
\lstset{language=kerml,
%    basicstyle=\tiny,
	keywordstyle={[1]\bfseries\color{blue}},  	%symbols
	keywordstyle={[2]\bfseries\color{blue}},  %BLESS keywords
	commentstyle={\itshape\color{commentcolor}},  %comments
	frame=single,
	frameround=tttt,
	backgroundcolor= \color{codebackgroundcolor},	
%	lineskip=-3pt,
%	escapechar=\#,  %use to insert some LaTeX command in the listing
%	numbers=left
	}

%switch font size to large temporarily for \lstinline
\newcommand{\lstk}[1]
{\lstset{language=kerml,basicstyle=\ttfamily\large}\lstinline{#1}\lstset{language=kerml,basicstyle=\ttfamily\listingsize}}

%define a language for Lean 4 listings (Lean4SFS/SFS.lean, Lean4SFS/DSL.lean)
\lstdefinelanguage{lean}
	{
	keywords={},
	morekeywords=[2]{
theorem, lemma, def, abbrev, structure, inductive, class, instance, axiom, opaque,
example, where, deriving, extends, mutual, namespace, section, end, open, import,
universe, variable, set_option, noncomputable, partial, unsafe, private, protected,
scoped, local, macro, macro_rules, syntax, notation, infix, infixl, infixr, prefix,
postfix, elab, elab_rules, attribute, by, do, let, fun, match, with, if, then, else,
have, show, suffices, from, calc, this, sorry, exact, apply, intro, intros, rfl, rw,
simp, constructor, cases, rcases, obtain, induction, refine, use, forall, Type, Prop,
Sort
},
	sensitive=true,
	morecomment=[l]{--},
	morecomment=[s]{/-}{-/},
	morestring=[b]",
	columns=[r]fixed,
	basicstyle=\ttfamily\listingsize,
	tabsize=2,
	}

\lstset{language=lean,
	keywordstyle={[2]\bfseries\color{blue}},
	commentstyle={\itshape\color{commentcolor}},
	frame=single,
	frameround=tttt,
	backgroundcolor= \color{codebackgroundcolor},
	}
%restore kerml as the document-wide default language -- the \lstset{language=lean,...} call
%above (needed to attach frame/commentstyle/etc. to the lean language) would otherwise silently
%make lean the default for every later \begin{lstlisting} in the book that doesn't say
%[language=kerml] explicitly. Select Lean listings per-call with \begin{lstlisting}[language=lean].
\lstset{language=kerml}


\renewcommand{\comment}[1]{{\color{commentcolor}{#1}}}

\newcommand{\ksec}[1]{\section{#1}}
%\newcommand{\ksec}[1]{\section{  [ #1 ]}\addcontentsline{toc}{section}{#1}}

%%%% LIST OF CODE %%%%
\newcommand\ecaption[1]{\addcontentsline{xmp}{example}{#1}}
\makeatletter
\newcommand\listofexamples{\chapter*{\huge{List of Example Code}}\@starttoc{xmp}}
\newcommand\l@example[2]{\@dottedtocline{1}{1.5em}{2.3em}{#1} {#2}}
\makeatother


\newcommand{\Finv}{\rotatebox[origin=C]{180}{$\mathsf{F}$}}

\newcommand\specialcaret{%
  \stackengine{0pt}{\ \,}{\scalebox{1.5}[2]{\raisebox{-1ex}{\string^}}}{O}{c}{F}{T}{L}}


%%%% LIST OF MATH %%%%
\newcommand\lom[1]{\addcontentsline{lom}{math}{#1}}
\makeatletter
\newcommand\listofmath{\chapter*{\huge{Metamath Metadefinitions and Metatheorems}}\@starttoc{lom}}
\newcommand\l@math[2]{\@dottedtocline{1}{1.5em}{2.3em}{#1} {#2}}
\makeatother

%%%% LIST OF INFERENCE RULES %%%%
\newcommand\loir[1]{\addcontentsline{loir}{rule}{#1}}
\makeatletter
\newcommand\listofrules{\chapter*{\huge{Semantic Definitions}}\@starttoc{loir}}
\newcommand\l@rule[2]{\@dottedtocline{1}{1.5em}{2.3em}{#1} {#2}}
\makeatother

%%% ToC APPEARANCE
\makeatletter
\renewcommand*{\l@subsection}{\@dottedtocline{2}{3.8em}{4em}} %\@dottedtocline{?seclevel ?}{?indent ?}{?numwidth ?}.
\renewcommand*{\l@section}{\@dottedtocline{1}{1.5em}{3.3em}} %\@dottedtocline{?seclevel ?}{?indent ?}{?numwidth ?}.
\makeatother
\renewcommand{\cftsecfont}{\rmfamily\mdseries\upshape}
\renewcommand{\cftsecpagefont}{\rmfamily\mdseries\upshape} % No bold!
\setlength{\cfttabnumwidth}{3em}  %make list of tables number width bigger
\setlength{\cftsubsecnumwidth}{20em}% Set length of number width in ToC for \subsection
\setlength{\cftsecnumwidth}{6em}% Set length of number width in ToC for \section
\setlength{\cftchapnumwidth}{2em}% Set length of number width in ToC for \chapter

%%% INDEX
\makeindex		% cause LaTeX to record index entries
\newcommand{\bemph}[1]		%BLESS emphasis
	{\emph{#1}\index{#1}}	%end of \bemph
\newcommand{\textidx}[1]  %text word into index
  {\texttt{#1}\index{#1@\texttt{#1}}}
%\newcommand{\remph}[1]		%Requirements emphasis
%	{\emph{#1}\index{#1}\footnote{requirement R\arabic{section}.\arabic{subsection}.\arabic{subsubsection}(\arabic{paragraph}):  \emph{#1}}%
%	\index{zzzR\arabic{section}.\arabic{subsection}.\arabic{subsubsection}(\arabic{paragraph})@R\arabic{section}.\arabic{subsection}.\arabic{subsubsection}(\arabic{paragraph}) #1}}	%end of \remph
%\newcommand{\rindex}[1]  %Requirements index	
%	{\index{#1}\footnote{requirement R\arabic{section}.\arabic{subsection}.\arabic{subsubsection}(\arabic{paragraph}):  \emph{#1}}%
%	\index{zzzR\arabic{section}.\arabic{subsection}.\arabic{subsubsection}(\arabic{paragraph})@R\arabic{section}.\arabic{subsection}.\arabic{subsubsection}(\arabic{paragraph}) #1}}	%end of \rindex

%THEOREM STYLES
\newtheoremstyle{semanticsstyle} %name of theorem style is "semanticsstyle"
	{3pt} %space-above
	{3pt} %space-below
	{\upshape} %body-style
	{-10pt} %indent
	{\scshape\raggedright} %head-style
	{ } %head-after-punct
	{ } %head-after-space
	{ } %head-full-spec

%\theoremstyle{semanticsstyle}
%\theoremstyle{plain}

% counter for verification conditions
%\newcounter{vccounter}
\declaretheoremstyle[
spaceabove=-10pt, spacebelow=-10pt,
headfont=\normalfont\bfseries,
notefont=\mdseries, notebraces={(}{)},
bodyfont=\normalfont,
postheadspace=2pt,
qed={}
]{blstyle}

\declaretheorem[name=Theorem,style=blstyle,numbered=no,
   shaded={rulecolor=green,rulewidth=2pt,bgcolor=white}]{theorem}
\declaretheorem[name={Lean Theorem},style=blstyle,numbered=no,
   shaded={rulecolor=green,rulewidth=2pt,bgcolor=white}]{leantheorem}
%\newtheorem{theorem}{Theorem} %[section]
\declaretheorem[name=Definition,style=blstyle,numbered=no,
   shaded={rulecolor=red,rulewidth=2pt,bgcolor=white}]{definition}
\declaretheorem[name={Lean Definition},style=blstyle,numbered=no,
   shaded={rulecolor=red,rulewidth=2pt,bgcolor=white}]{leandefinition}
\declaretheorem[name=Metadefinition,style=blstyle,numbered=no,
   shaded={rulecolor=blue!50!black,rulewidth=2pt,bgcolor=white}]{metadefinition}
%\newtheorem{definition}{Definition} %[section]
\declaretheorem[name=Metatheorem,style=blstyle,numbered=no,
  shaded={rulecolor=orange,rulewidth=2pt,bgcolor=white},]{lemma}
%\newtheorem{lemma}{Lemma}  %[theorem]{Lemma}
\declaretheorem[name=Axiom,style=blstyle, numbered=no,
   shaded={rulecolor=yellow!50!black,rulewidth=2pt,bgcolor=white}]{axiom}
\declaretheorem[name={Lean Axiom},style=blstyle, numbered=no,
   shaded={rulecolor=yellow!50!black,rulewidth=2pt,bgcolor=white}]{leanaxiom}
%\newtheorem{axiom}{Axiom} %[section]
\declaretheorem[name={Verification Condition},style=blstyle]{vc}
%\newtheorem{vc}{Verification Condition} %[vccounter]

% shorter align environment replacement
\newenvironment{salign}
  {\vspace{-8pt}
  \[
  }
  {
  \]
  \vspace{-10pt}
  }

\newcommand{\leanTitle}[2]{\emph{\large\tt\color{leancolor}#1}~~{#2}
  \index{zzzl@\color{leancolor}Lean!\tt#1}\lom{\textcolor{leancolor}{\tt#1}~#2}}

%replace "\label{def:" or "\label{the:" with  "\blTitle{" on all metamath-generated lemmas (BLESS theorems)
\newcommand{\blTitle}[2]{\emph{\large\tt\color{blcolor}#1}~~{#2}
  \index{zzzbl@\color{blcolor}Metatheorems!\tt#1}\lom{\textcolor{blcolor}{\tt#1}~#2}}

%replace "\label{def:" with  "\otherTitle{" on all metamath-generated lemmas (everything else)
\newcommand{\otherTitle}[2]{\emph{\large\tt\color{setmmcolor}#1}~~{#2}
  \index{zzzmm@\color{setmmcolor}Metatheorems (from set.mm)!\tt#1}\lom{\textcolor{setmmcolor}{\tt#1}~#2}}
   
\newcommand{\ruleTitle}[2]{~{\bf\large #1}~~{\it#2}
  \index{zzzrule@Semantic Definitions!#1}\loir{{\bf#1}~{\it#2}}}
   
%replace "\ref{eq:" with "\mmref{"
\newcommand{\mmref}[1]{{\color{orange}\ref{#1}}}   

% replace "\label{eq:"  with "\label{"

% replace "\notag \\ && & \qquad " with " "

%put Metamath theorem identifier in index, and make its label
\newcommand{\mmlabel}[1]
  {\label{eq:#1}\index{Metatheorems!#1}
}% end of \mmlabel 

% GLOSSARY ENTRIES
\setlength{\fboxrule}{2pt}
\newcommand{\gloss}[2]
  {\index{#1|textbf} {\color{glossarycolor}\begin{center}\fbox{\parbox{\linewidth-1in}{\textbf{#1:}~~#2}}\end{center}}}

%PREVENT RESET OF PARAGRAPH COUNTER
%\usepackage{chngcntr}
\counterwithout{paragraph}{section}
\counterwithout{paragraph}{subsection}
\counterwithout{paragraph}{subsubsection}

%COUNTER RESET
\makeatletter
%\@addtoreset{paragraph}{chapter}
\@addtoreset{paragraph}{section}
%\@addtoreset{legalityrule}{chapter}
%\@addtoreset{legalityrule}{section}
%\@addtoreset{namingrule}{chapter}
%\@addtoreset{namingrule}{section}
%\@addtoreset{consistencyrule}{chapter}
%\@addtoreset{consistencyrule}{section}
%\@addtoreset{semantic}{chapter}
%\@addtoreset{semantic}{section}
%\@addtoreset{chapter}{part}
\makeatother

%PARAGRAPH NUMBERS
% pn
\newlength{\pnwidth}	%declare paragraph number width
\newcommand\pn{
	\stepcounter{paragraph}
	\settowidth{\pnwidth}{(\arabic{paragraph})~~}  %get width of (1)
	\hspace{-\pnwidth}
	\raisebox{1pt}
	{\color{paragraphnumbercolor}\small(\arabic{paragraph})}~}
%\newcommand\pn{}

%SYMBOLS FOR << AND >>
\newcommand{\llx}{\ll\!\!}	%left <<
\newcommand{\xgg}{\!\!\gg} 	%right >>

%WORDS THAT ARE MATH OPERATORS
\DeclareMathOperator*{\w}{\mathsf{wp}}
\DeclareMathOperator*{\biff}{\lrarrow}
\DeclareMathOperator*{\interferencefree}{interference-free}
%\DeclareMathOperator*{\bif}{if}
%\DeclareMathOperator*{\bfi}{fi}
\DeclareMathOperator*{\bguard}{)~>}
\DeclareMathOperator*{\band}{\wedge}
\DeclareMathOperator*{\bor}{\vee}
%\DeclareMathOperator*{\bxor}{xor}
%\DeclareMathOperator*{\bnot}{not}
\DeclareMathOperator*{\modulus}{\mathsf{mod}}
\DeclareMathOperator*{\bdiv}{\mathsf{div}}
\DeclareMathOperator*{\bcard}{\mathsf{card}}
\DeclareMathOperator*{\bmin}{\mathsf{min}}
\DeclareMathOperator*{\brem}{\mathsf{rem}}
\DeclareMathOperator*{\intervalstart}{\mathsf{start}}
\DeclareMathOperator*{\intervalend}{\mathsf{end}}
\DeclareMathOperator*{\plusarrow}{+=>}
\DeclareMathOperator*{\btops}{\mathsf{tops}}
\DeclareMathOperator*{\bnow}{\mathsf{now}}
\DeclareMathOperator*{\dotdot}{..}
\DeclareMathOperator*{\commadot}{,.}
\DeclareMathOperator*{\dotcomma}{..\,}
\DeclareMathOperator*{\commacomma}{,,}

% LATTICE DRAWING
% counters for lattice drawing
\newcounter{xnode}		%node center x coordinate
\newcounter{ynode}		%node center y coordinate
\newcounter{lwidth}		%lattice width
\newcounter{lheight}		%lattice height
\newcounter{noderow}	%node row
\newcounter{numrows} %number of rows
\newcounter{node}			%which node in row
\newcounter{numnodesinrow}	%number of nodes in row
\newcounter{nodespacing}	%length between node centers
\newcounter{circlediameter}	%diameter of circles drawn
\newcommand{\nodetextscalefactor}{1}	%scale factor for scalebox of text in lattice nodes

% DRAW NODE
\newcommand{\drawnode}[4] %{row}{numrows}{node}{numnodesinrow}
	{
	\setcounter{xnode}{#3 * (\value{lwidth} / (#4 + 1))}
	\setcounter{ynode}{\value{lheight} - #1 * (\value{lheight} / (#2 + 1))}
	\put(\value{xnode},\value{ynode}){\circle*{\value{circlediameter}}}
	{\color{black} 
		\put(\value{xnode},\value{ynode}){\circle{\value{circlediameter}}}
		}
	}
	%end of \drawnode

% DRAW NODE WITH TEXT
\newcommand{\drawtextnode}[5] %{row}{numrows}{node}{numnodesinrow}{text}
	{
	\setcounter{xnode}{#3 * (\value{lwidth} / (#4 + 1))}
	\setcounter{ynode}{\value{lheight} - #1 * (\value{lheight} / (#2 + 1))}
	\put(\value{xnode},\value{ynode}){\circle*{\value{circlediameter}}}
	{\color{black} 
		\put(\value{xnode},\value{ynode}){\circle{\value{circlediameter}}}
		%move 1/2 width left and 1/2 height down
%		\addtocounter{xnode}{0-(\value{circlediameter}/4)}
%		\addtocounter{ynode}{0-(\value{circlediameter}/4)}
		\put(\value{xnode},\value{ynode}){#5}  %\scalebox{\nodetextscalefactor}{#5}}
		}
	}
	%end of \drawnode

% DRAW NODE ROW
\newcommand{\drawnoderow}[3] %{row}{numrows}{numnodesinrow}
	{
	\setcounter{xnode}{\value{lwidth} / (#3 + 1)}
	\setcounter{ynode}{\value{lheight} - #1 * (\value{lheight} / (#2 + 1))}
	\multiput(\value{xnode},\value{ynode})(\value{xnode},0){#3}{\circle*{\value{circlediameter}}}
	{\color{black} 
		\multiput(\value{xnode},\value{ynode})(\value{xnode},0){#3}{\circle{\value{circlediameter}}}
		}
	}
	%end of \drawnode

% DRAW ARC TO LOWER ROW
	\newcounter{xsource}	%x coordinate of source node
	\newcounter{ysource}	%y coordinate of source node
	\newcounter{xmid}	%x coordinate of midpoint node
	\newcounter{ymid}	%y coordinate of midpoint node
	\newcounter{xdest}	%x coordinate of destination node
	\newcounter{ydest}	%y coordinate of destination node
\newcommand{\drawarc}[6]	%{row}{numrows}{fromnode}{nodesinfromrow}{tonode}{nodesintorow}
	{
	\setcounter{xsource}{#3 * \value{lwidth} / (#4 + 1)}	%fromnode*width/(nodesinfromrow+1)
	\setcounter{xdest}{#5 * \value{lwidth} / (#6 + 1)}	%tonode*width/(nodesintorow+1)
	\setcounter{xmid}{(\value{xsource}+ \value{xdest})/2}					%midpoint is average
	\setcounter{ysource}{\value{lheight} - #1 * (\value{lheight} / (#2 + 1))
				 %- (\value{circlediameter}/2)
				 }	%height-row*height/(numrows+1)
	\setcounter{ydest}{\value{lheight} - (#1 + 1) * (\value{lheight} / (#2 + 1))
				 %+ (\value{circlediameter}/2)
				 }	%height-(row+1)*height/(numrows+1)
	\setcounter{ymid}{(\value{ysource}+ \value{ydest})/2}					%midpoint is average
	\ifthenelse{#4 > #6}			% nodesinfromrow < 	nodesintorow		 
		{\setcounter{ymid}{(\value{ysource} + (3*\value{ydest}))/4}	}				
		{\setcounter{ymid}{(\value{ydest} + (3*\value{ysource}))/4}		}				
%	\qbezier(\value{xsource},\value{ysource})
%			(\value{xmid},\value{ysource})
%			(\value{xmid}, \value{ymid})
%	\qbezier(\value{xmid}, \value{ymid})
%			(\value{xmid}, \value{ydest})
%			(\value{xdest}, \value{ydest})	
	\qbezier(\value{xsource}, \value{ysource})
			(\value{xmid}, \value{ymid})
			(\value{xdest}, \value{ydest})	
	}
	%end of \drawarc

% MATHEMATICAL MEANING DEFINITION "M" (must be in math mode)
%make interval shape
\newcommand{\interval}[1]{\ensuremath{
	{\color{red}\underline{\mathbf{\mathfrak{#1}}}}}}		
%make interval shape with subscript
\newcommand{\intervalsub}[2]{\ensuremath{
	{\color{red}\underline{\mathbf{\mathfrak{#1}}}_{#2}}}}
%make M[[]]		
\newcommand{\md}[1]{\ensuremath{\llbracket\mathfrak{M}\models~#1\rrbracket}}	
%make Mi[[]]
\newcommand{\mdi}[1]{\ensuremath{\llbracket\mathfrak{M}_{\interval{i}}\models~#1\rrbracket}}	
%make Ms[[]]
\newcommand{\mds}[1]{\ensuremath{\llbracket\mathfrak{M}_{s}\models~#1\rrbracket}}	
%make Mx[[]]
\newcommand{\mdx}[2]{\ensuremath{\llbracket\mathfrak{M}_{#2}\models~#1\rrbracket}}	

%SEMANTICS COMMENTS
\newcommand{\ssc}[1]{{\color{commentcolor}\textit{(#1)}}}

%these eliminate unnecessary stuff from BAv2
\newcommand{\JP}[1]{}
\newcommand{\quoteBA}[2]{#2}
\newcommand{\refBA}[1]{}
\newcommand{\twoBA}[2]{}
\newcommand{\notBA}[1]{}
\newcommand{\reconciliation}[1]{} 

%create mini tables of contents for each part (minitoc)
%\doparttoc
\newcommand{\riota}{\j}

%BNF environment for grammar text
\newenvironment{bnf}
	{\VerbatimEnvironment\begin{BVerbatim}[commandchars=\\\{\},fontfamily=courier,obeytabs=true,tabsize=3]}
	{\end{BVerbatim}}

%put tilde in the middle
\newcommand{\midtilde}{\raisebox{-0.3\baselineskip}{\textbf\textasciitilde}}

